\chapter{Subquadratic Sparse Attention (SSA)}

\section{Public framing}

Subquadratic markets \emph{Subquadratic Sparse Attention} (SSA) as an attention mechanism whose cost scales
favorably with context length by avoiding fully dense all-pairs attention \cite{subq2026ssa}.
Their technical article states that SSA uses \emph{content-dependent selection}: for each query, the model
chooses which positions receive attention, then applies attention on that subset rather than on every pair.
They claim exact attention on selected positions (not a softmax surrogate), with proprietary selection rules.

\section{A tradeoff taxonomy}

Efficient long-context modeling faces a recurring tension:
\begin{itemize}
  \item \textbf{Fixed-pattern sparse masks} (sliding windows, strided patterns): cheap, but routing is chosen
    by position in advance; relevant tokens outside the pattern may be unreachable.
  \item \textbf{Recurrent / compressed states}: linear-time updates, but finite state may lose verbatim access
    to distant tokens.
  \item \textbf{Hybrid stacks}: combine efficient layers with occasional dense attention; dense layers may
    still dominate cost at extreme $n$.
\end{itemize}

\section{Reported throughput and benchmarks (vendor-reported)}

Subquadratic publishes wall-clock prefill speedups versus FlashAttention-2 on NVIDIA B200 GPUs at multiple
context lengths, e.g.\ $7.2\times$ at 128K through $52.2\times$ at 1M tokens \cite{subq2026ssa}. They also report
accuracy figures on long-context retrieval benchmarks (e.g.\ RULER, MRCR v2) alongside coding benchmarks (e.g.\
SWE-Bench Verified). Readers should treat headline tables as vendor-reported until independent reproductions and
full model cards are available.

\paragraph{Caveat.}
Media coverage has highlighted demand for independent verification of extreme efficiency claims
\cite{venturebeat2026subq}. This document is educational material aligned with the LLMs repository's research
notes; it is not an endorsement of every numerical claim.
